\documentclass{article}
\usepackage[margin=1in]{geometry}
\usepackage{graphicx} % Required for inserting images
\usepackage{booktabs}
\usepackage{float}
\usepackage{amsmath}

% Figures: upload auto_mpg_transformed_regression.png and concrete_transformed_regression.png
% (from project2/results/plots/) to the same Overleaf folder as this file.

\title{Transformed Regression}
\author{Adithya Lakshmikanth}
\date{September 2026}

\begin{document}

\maketitle

\section{Introduction}

Transformed regression keeps the linear model and its efficient least-squares training but applies a transformation $f$ to the response:
\[
f(y) = \mathbf{b}\cdot\mathbf{x} + \epsilon, \qquad \hat{y} = f^{-1}(\mathbf{b}\cdot\mathbf{x}).
\]
It was fitted with ScalaTion's \texttt{TranRegression} class on two of the three datasets, Auto MPG and Concrete, using all available predictors plus an intercept. For each dataset, five models were compared: plain \texttt{Regression} (no transformation) and \texttt{TranRegression} with a log, square-root, reciprocal, and Box-Cox transformation,
\[
f_\lambda(y) = \frac{y^{\lambda}-1}{\lambda}\ (\lambda \neq 0), \qquad f_0(y) = \log y .
\]
The transformed variable is the response only: fuel economy (\texttt{mpg}, in MPG) for Auto MPG and compressive strength (\texttt{compressive\_strength\_mpa}, in MPa) for Concrete. Predictor variables are left untransformed. Box--Cox was selected as the required alternative to Yeo--Johnson; since every observed response is strictly positive, Yeo--Johnson's extension to nonpositive values is not needed. The reciprocal transformation is included as an additional comparison.

Because $\lambda = 1$ is a shift of $y$, the Box-Cox family contains ``no transformation'' as a special case. The Box-Cox parameter was tuned over $\lambda \in \{-1, -0.75, \dots, 2\}$ by 5-fold cross-validation on the training set only. Predictions were mapped back with $f^{-1}$, so the comparison tables report $R^2$, RMSE, and MAE on the \emph{original} response scale. Full-data fits are descriptive; held-out results come from separate models fitted only to the training rows. Both datasets use an 80/20 split generated by Scala's random shuffle with seed 42. For Auto MPG, this reproduces the split used by the existing ScalaTion Ridge and Lasso code. Auto MPG uses 314 training and 78 test rows; Concrete uses 824 training and 206 test rows. The already shuffled training rows are divided into five contiguous folds, and CV RMSE pools squared errors across all validation rows. All observed responses are positive. For $\lambda\ne0$, a valid Box--Cox inverse requires $1+\lambda z>0$, even when an integer inverse power would return a finite number outside this domain. Square-root and positive-response reciprocal models require $z\geq0$ and $z>0$, respectively. Candidates with any invalid or nonfinite validation prediction are excluded. An invalid full-data or test prediction makes that entire model/split score unavailable; rows are never silently discarded.

Auto MPG and Concrete satisfy the requirement to apply transformed regression to two of the three datasets. They provide a useful comparison because the observed benefit is substantial for Auto MPG and small for Concrete. No transformed-regression result is claimed for Airfoil.

The predictions use direct inverse transformation, without a retransformation bias correction. In general, $f^{-1}(E[f(y)\mid x])$ is not $E[y\mid x]$; these fits minimize squared error on the transformed scale, not the original scale. Original-scale validation measures the practical prediction error of this choice. The chosen $\lambda$ minimizes training CV RMSE; the held-out test set is not used for tuning.

All predictions, metrics, and CV scores passed independent NumPy verification.

\section{Auto MPG}

\begin{table}[H]
\centering
\small
\begin{tabular}{lrrr|rrr}
\toprule
 & \multicolumn{3}{c|}{In-sample ($n=392$)} & \multicolumn{3}{c}{Held-out test} \\
Model & $R^2$ & RMSE & MAE & $R^2$ & RMSE & MAE \\
\midrule
Regression (no transform)      & 0.8215 & 3.294 & 2.499 & 0.7835 & 3.528 & 2.631 \\
TranRegression, log            & 0.8562 & 2.956 & 2.102 & 0.8216 & 3.202 & 2.285 \\
TranRegression, sqrt           & 0.8437 & 3.082 & 2.252 & 0.8059 & 3.340 & 2.426 \\
TranRegression, reciprocal     & 0.8490 & 3.029 & 2.137 & 0.8069 & 3.331 & 2.274 \\
TranRegression, Box-Cox $\lambda=-0.5$ & \textbf{0.8612} & \textbf{2.904} & \textbf{2.043} & \textbf{0.8279} & \textbf{3.145} & \textbf{2.216} \\
\bottomrule
\end{tabular}
\caption{Auto MPG: original-scale quality of fit for plain and transformed regression (ScalaTion).}
\label{tab:tran_autompg}
\end{table}

Every transformation improved on plain regression, both in-sample and on the test set. Cross-validation selected $\lambda = -0.5$ (CV RMSE 2.995 versus 3.391 at $\lambda = 1$), and this model reduced test RMSE from 3.528 to 3.145~MPG, a 10.9\% reduction. The selected $\lambda$ lies between the log ($\lambda = 0$) and reciprocal ($\lambda = -1$) transformations. The reciprocal alternative models fuel consumption (gallons per mile, $1/\text{MPG}$) directly; its improvement over plain regression supports examining nearby powers.

Figure~\ref{fig:tran_autompg} compares full-data fits and residuals. The Box--Cox model improves the fit, but these plots are descriptive and do not replace the held-out evaluation. On the transformed scale, the selected full-data model has $R^2=0.890$. Model year ($t=15.9$) and weight ($t=-10.7$) have the largest absolute $t$ statistics among its predictors. Acceleration has $p=0.23$. These are conventional OLS statistics conditional on the selected transformation, without adjustment for choosing $\lambda$ from the data; they are descriptive rather than confirmatory. Coefficients and their standard errors are computed on the transformed scale and are not in MPG units. Origin is treated as a numeric predictor, matching the existing regression setup.

\begin{figure}[H]
\centering
\includegraphics[width=\textwidth]{auto_mpg_transformed_regression.png}
\caption{Auto MPG: Box-Cox $\lambda$ tuning curve (left), actual vs.\ predicted MPG (center), and residuals vs.\ fitted values (right) for plain regression and the Box-Cox $\lambda=-0.5$ model.}
\label{fig:tran_autompg}
\end{figure}

\section{Concrete}

\begin{table}[H]
\centering
\small
\begin{tabular}{lrrr|rrr}
\toprule
 & \multicolumn{3}{c|}{In-sample ($n=1030$)} & \multicolumn{3}{c}{Held-out test} \\
Model & $R^2$ & RMSE & MAE & $R^2$ & RMSE & MAE \\
\midrule
Regression (no transform)      & 0.6155 & 10.354 & 8.215 & 0.5957 & 10.311 & 8.125 \\
TranRegression, log            & 0.4559 & 12.317 & 9.442 & 0.4740 & 11.761 & 9.305 \\
TranRegression, sqrt           & 0.5800 & 10.822 & 8.637 & 0.5582 & 10.779 & 8.571 \\
TranRegression, reciprocal     & --- & --- & --- & --- & --- & --- \\
TranRegression, Box-Cox $\lambda=1.25$ & \textbf{0.6218} & \textbf{10.269} & \textbf{8.074} & \textbf{0.6055} & \textbf{10.185} & \textbf{7.970} \\
\bottomrule
\end{tabular}
\caption{Concrete: original-scale quality of fit. Reciprocal scores are unavailable because the positive-response inverse domain fails for 21 of 1,030 full-data predictions and 1 of 206 test predictions.}
\label{tab:tran_concrete}
\end{table}

Concrete benefits much less from response transformation. Log and square-root transformations increased prediction error relative to plain regression: log reduced test $R^2$ from 0.596 to 0.474, and square root reduced it to 0.558. The reciprocal transformation failed badly. The model on $1/y$ produces some nonpositive fitted values, outside the inverse domain for positive strengths. Near-zero positive values can also produce very large predictions. This makes the reciprocal model unsuitable here. Cross-validation chose $\lambda = 1.25$, slightly \emph{above} 1, which expands rather than compresses the response. The resulting gain is marginal: test RMSE fell only from 10.311 to 10.185~MPa (1.2\%). The valid portion of the CV curve in Figure~\ref{fig:tran_concrete} is nearly flat between $\lambda=0.75$ and $1.5$ (CV RMSE 10.43 at the optimum versus 10.51 with no transformation), so the data provide little evidence for any response transformation.

These results show that transforming only the response offers limited benefit for this linear predictor specification. They do not establish the source of any remaining nonlinearity. Transforming predictors or adding interactions could be investigated separately through symbolic regression. The small improvement on one held-out split is not evidence of a statistically significant gain.

\begin{figure}[H]
\centering
\includegraphics[width=\textwidth]{concrete_transformed_regression.png}
\caption{Concrete: Box-Cox $\lambda$ tuning curve (left; invalid candidates omitted), actual vs.\ predicted strength (center), and residuals vs.\ fitted values (right) for plain regression and the Box-Cox $\lambda=1.25$ model.}
\label{fig:tran_concrete}
\end{figure}

\section{Summary}

For Auto MPG, training CV selected Box--Cox $\lambda=-0.5$, and test RMSE fell by 10.9\% relative to plain regression. For Concrete, $\lambda=1.25$ reduced test RMSE by only 1.2\%; log and square-root transformations worsened test error, and the reciprocal model produced invalid predictions. These conclusions apply to the tested models and this fixed split, not to all possible ScalaTion or symbolic-regression models. A response transformation should be chosen using validation in the units of the prediction task.

\end{document}
